% =====================================================================
%  CHƯƠNG 8 — TRIỂN KHAI THIẾT BỊ BIÊN VÀ ĐÁNH GIÁ
%  Phụ trách: Nhật (tối ưu mô hình) + Đức (đóng gói, worker)
%  Viết high-level, không nhắc mã nguồn/hàm/đường dẫn.
%  INT8: CHƯA thực hiện — chỉ nêu ở hướng phát triển (Chương 9).
%  Số Pi5 đo trực tiếp ngày 2026-09-13.
% =====================================================================
\chapter{Triển Khai Trên Thiết Bị Biên và Đánh Giá}
\label{ch:trienkhai}

\section{Nền Tảng Phần Cứng Mục Tiêu}
\label{sec:trienkhai-phancung}

Raspberry Pi 5 được chọn vì CPU bốn nhân Cortex-A76 đủ chạy YOLO cỡ nano qua ONNX Runtime mà
không cần GPU rời (Mục~\ref{sec:pi5bench}), chi phí thấp và có sẵn giao tiếp camera CSI lẫn
GPIO. Mọi mô hình được huấn luyện trên máy tính có GPU; các phép đo hiệu năng chạy
trên CPU của cả hai máy để so cùng điều kiện (Bảng~\ref{tab:moi-truong}).

\begin{table}[htbp]
\centering
\caption{Cấu hình phần cứng và phần mềm dùng trong thực nghiệm.}
\label{tab:moi-truong}

\vspace{4pt}
\begin{tabular}{lll}
\toprule
\textbf{Thành phần} & \textbf{Máy tính} & \textbf{Raspberry Pi 5} \\
\midrule
CPU          & Intel Core i7-13700K (16 nhân, 24 luồng) & 4 nhân Cortex-A76 2,4 GHz \\
Bộ nhớ       & 64 GB & 16 GB \\
GPU          & NVIDIA RTX 4070 (chỉ khi huấn luyện) & Không \\
Hệ điều hành & Windows 11 Pro & Debian GNU/Linux 13, kernel 6.18 \\
ONNX Runtime & 1.28 & 1.30 \\
Backend      & FastAPI, PostgreSQL & FastAPI, SQLite \\
\bottomrule
\end{tabular}
\end{table}

\section{Tối Ưu và Đóng Gói Mô Hình}
\label{sec:trienkhai-toiuu}

Ba mô hình đọc ký tự, phân loại loại xe và phân loại kiểu dáng chạy dưới ONNX Runtime, không
phụ thuộc PyTorch khi suy luận; các detector biển số và định vị xe chạy qua Ultralytics. Bản
ONNX cho dự đoán trùng khớp bản huấn luyện (Mục~\ref{sec:ocr-charset37}) và chạy được trên CPU
ARM, macOS lẫn Windows, nên cùng một mô hình dùng chung cho máy tính và thiết bị biên mà không
biên dịch lại. Các biện pháp tối ưu đã áp dụng gồm xuất ONNX, đặt chung số luồng cho mọi mô
hình và dùng MobileNetV3-Small cho bước kiểu dáng; lượng tử hóa INT8 chưa thực hiện.

Worker tại cổng gồm bốn thành phần: nhận tín hiệu có xe (nút bấm hoặc cảm biến qua GPIO), đọc
khung ảnh từ camera, chạy suy luận tại chỗ, và gửi kết quả kèm ảnh về backend, xác thực bằng
khóa riêng của thiết bị biên. Backend nhận kết quả này không suy luận lại, chỉ ghi nhận và
tạo phiên, bảo đảm kết quả đến từ đúng mô hình trên thiết bị biên; mỗi lỗi chỉ bỏ một lượt
chụp mà không làm dừng worker.

\section{Đánh Giá Theo Thành Phần}
\label{sec:trienkhai-thanhphan}

Bảng~\ref{tab:eval-components} tổng hợp kết quả từng thành phần trên các tập kiểm tra riêng
đã trình bày ở Chương~\ref{ch:phathien} đến~\ref{ch:phanloai}.

\begin{table}[htbp]
\centering
\caption{Tổng hợp hiệu năng các thành phần của pipeline nhận diện phương tiện.}

\label{tab:eval-components}
\begin{tabular}{llll}
\toprule
\textbf{Thành phần} & \textbf{Mô hình} & \textbf{Chỉ số} & \textbf{Kết quả} \\
\midrule
Phát hiện biển số & YOLOv8n (A1, TB 3 seed)
  & mAP@0,5        & 0,989 \\
  & & mAP@0,5:0,95   & 0,864 \\
\midrule
OCR biển số & CRNN 37 ký tự (có ``Đ'')
  & Accuracy (a1, 172 biển)      & 91,9\% \\
  & & Accuracy (vn\_plate, 96 biển) & 80,2\% \\
  & & Accuracy (topkek, 666 biển)  & 56,6\% \\
\midrule
Phân loại loại xe & ResNet18
  & Accuracy       & 98,66\% \\
  & & F1-macro       & 96,19\% \\
\midrule
Phân loại kiểu dáng & MobileNetV3-Small (B5, 3 lớp)
  & Accuracy       & 90,07\% \\
  & & F1-macro       & 90,69\% \\
\midrule
Màu nền biển số & HSV + ngưỡng màu
  & Accuracy       & 100\% (trắng/đỏ, 178 crop) \\
\bottomrule
\end{tabular}
\end{table}

Với detector, con số đáng tin hơn là mAP@0,5 bằng \textbf{0,983 sau khử trùng lặp gần}
(Bảng~\ref{tab:det-dedup}); với OCR, con số sát điều kiện triển khai nhất là
\textbf{91,9\% trên a1} (CER 0,015).

\section{Đánh Giá End-to-End trên Raspberry Pi 5}
\label{sec:trienkhai-e2e}

Toàn luồng được đo trực tiếp trên Raspberry Pi~5 với 50 ảnh camera cổng thật, gửi qua API
thật. Độ trễ đầu cuối trung vị \textbf{463,4 ms} (Bảng~\ref{tab:e2e-api}), trong đó phần
ngoài suy luận (HTTP, mã hóa JPEG, ghi cơ sở dữ liệu) chỉ 17,5 ms, nên muốn nhanh hơn phải
giảm thời gian mô hình chứ không phải tối ưu backend. Lượt đầu tốn 4,8 s do nạp mô hình, các
lượt sau ổn định.

\begin{table}[htbp]
\centering
\caption{Độ trễ end-to-end mỗi lượt xe trên Raspberry Pi~5 (50 ảnh cổng thật).}

\label{tab:e2e-api}
\begin{tabular}{lr}
\toprule
\textbf{Chỉ số} & \textbf{Giá trị} \\
\midrule
Độ trễ API đầu cuối (trung vị)      & 463,4 ms \\
Độ trễ API đầu cuối (p90)           & 487,9 ms \\
Trong đó: suy luận backend           & 445,8 ms \\
Trong đó: HTTP, JPEG, ghi cơ sở dữ liệu & 17,5 ms \\
Lượt đầu tiên (nạp mô hình)          & 4812,9 ms \\
Suy luận trên 50 ảnh (trung vị)      & 464,4 ms \\
Điện năng mỗi lượt (toàn phần)       & 3,121 J \\
Điện năng mỗi lượt (phần tăng thêm)  & 2,388 J \\
\bottomrule
\end{tabular}
\end{table}

Trên cả 50 ảnh, kết quả trên Pi~5 khớp tuyệt đối máy tính ở cả ba đầu ra: chuỗi biển số,
loại xe và kiểu dáng (50/50). Phép đo này chưa phải accuracy toàn luồng vì 50 ảnh chưa có nhãn
chuỗi ký tự chuẩn.

\section{Hiệu Năng Suy Luận Trên Thiết Bị Biên}
\label{sec:trienkhai-hieunang}

Mỗi cấu hình được đo sau khi SoC nguội về 58°C, lấy trung vị 30 lượt sau 5 lượt làm nóng,
trên cùng ảnh mẫu 12 biển như khi đo trên máy tính.

Pi~5 có bốn lõi nên tăng từ một lên bốn luồng mỗi mô hình giảm mạnh độ trễ các mô hình nặng
(Bảng~\ref{tab:pi-model}), trái với máy dev nơi số luồng cao gây tranh chấp. Toàn pipeline ở
bốn luồng nhanh nhất là 495,3 ms; ngay cả cấu hình một luồng an toàn nhất (813,9 ms) cũng
dưới ngân sách hai giây (Bảng~\ref{tab:pi-pipeline}). Hai mô hình YOLO (phát hiện biển và
định vị xe) chi phối tổng thời gian, nên đổi bước kiểu dáng sang MobileNetV3-Small giảm riêng
bước đó nhưng ít ảnh hưởng tổng.

\begin{table}[htbp]
\centering
\caption{Độ trễ từng mô hình ONNX chạy riêng trên Raspberry Pi~5 (trung vị, ms).}

\label{tab:pi-model}
\begin{tabular}{lrrr}
\toprule
\textbf{Model} & \textbf{Kích thước} & \textbf{1 luồng} & \textbf{4 luồng} \\
\midrule
Phát hiện biển (YOLOv8n)      & 12,2 MB & 360,9 & 142,9 \\
OCR biển số (CRNN)            &  1,7 MB &   6,1 &   2,3 \\
Loại xe (ResNet18)           & 44,7 MB & 129,0 &  45,3 \\
Kiểu dáng (ResNet18)         & 44,7 MB & 129,6 &  44,8 \\
Kiểu dáng (MobileNetV3-Small) &  6,1 MB &  12,5 &   5,8 \\
\bottomrule
\end{tabular}
\end{table}

\begin{table}[htbp]
\centering
\caption{Độ trễ toàn pipeline trên Raspberry Pi~5 theo cấu hình và số luồng (trung vị, ms, ảnh 12 biển).}

\label{tab:pi-pipeline}
\begin{tabular}{lrrr}
\toprule
\textbf{Cấu hình} & \textbf{1 luồng} & \textbf{2 luồng} & \textbf{4 luồng} \\
\midrule
Định vị bản Ultralytics, kiểu dáng ONNX   & 813,9 & 526,3 & 495,3 \\
ONNX thuần, kiểu dáng ResNet18            & 1067,4 & 611,6 & 535,0 \\
ONNX thuần, kiểu dáng MobileNetV3-Small   & 1082,4 & 610,9 & 550,4 \\
\bottomrule
\end{tabular}
\end{table}

\textbf{RAM, công suất và nhiệt độ.} Ở cấu hình bốn luồng, đỉnh RAM của tiến trình khoảng 720
đến 785 MB, thoải mái trong 16 GB. Công suất trung vị tăng từ khoảng 3,6 W (một luồng) lên
khoảng 6,1 W (bốn luồng), so với khoảng 1,5 W lúc nghỉ. Nhiệt độ SoC đạt tối đa khoảng 75°C
khi chạy benchmark và khoảng 83°C khi tải liên tục trên ảnh thật; suốt phép đo không ghi nhận
hạ xung hay sụt áp, xác nhận tản nhiệt chủ động đủ giữ hiệu năng ổn định. Bản
MobileNetV3-Small chạy mát hơn (3 mẫu chạm ngưỡng nhiệt mềm so với 13 của ResNet18), một lợi
thế cho vận hành liên tục.

\section{So Sánh PC và Thiết Bị Biên}
\label{sec:trienkhai-sosanh-pc-bien}

Trên cùng 50 ảnh camera cổng thật, cùng bốn luồng và cùng cấu hình triển khai, Pi~5 chậm hơn
máy tính khoảng 9,7 lần theo trung vị (464,4 ms so với 48,1 ms với MobileNetV3-Small), phản
ánh chênh lệch năng lực CPU; cả hai đều nằm dưới ngân sách hai giây (Bảng~\ref{tab:pc-vs-pi}).
Đổi bước kiểu dáng sang MobileNetV3-Small giữ nguyên trung vị nhưng hạ phân vị 90 trên Pi~5
(487,7 ms so với 539,1 ms).

\begin{table}[htbp]
\centering
\caption{So sánh độ trễ một lượt xe trên 50 ảnh camera cổng thật (trung vị và phân vị 90, ms).}

\label{tab:pc-vs-pi}
\begin{tabular}{llrr}
\toprule
\textbf{Nền tảng} & \textbf{Mô hình kiểu dáng} & \textbf{Trung vị} & \textbf{Phân vị 90} \\
\midrule
Raspberry Pi~5        & MobileNetV3-Small & 464,4 & 487,7 \\
Raspberry Pi~5        & ResNet18          & 463,0 & 539,1 \\
Máy tính (i7-13700K)  & MobileNetV3-Small &  48,1 &  54,0 \\
Máy tính (i7-13700K)  & ResNet18          &  49,0 &  59,4 \\
\bottomrule
\end{tabular}
\end{table}

\section{So Sánh Với Công Trình Liên Quan}
\label{sec:trienkhai-sosanh-lienquan}

Bảng~\ref{tab:compare-related} đối chiếu đề tài với ba công trình đã khảo sát ở
Chương~\ref{ch:lienquan}.

\begin{table}[htbp]
\centering
\caption{So sánh đề tài với các công trình liên quan về ALPR và bãi giữ xe.}
\label{tab:compare-related}

{\footnotesize Nguồn: số các công trình liên quan từ \autocite{tran2025vietnamlpr,pradhan2025iotparking,ammar2023multistage}; số của đề tài từ kết quả thực nghiệm.}
\begin{tabular}{p{3cm}p{2.5cm}p{2.5cm}p{3.5cm}}
\toprule
\textbf{Công trình} & \textbf{Nền tảng} & \textbf{Độ chính xác biển} & \textbf{Phạm vi hệ thống} \\
\midrule
Tran và Bui~\autocite{tran2025vietnamlpr}
  & Raspberry Pi 4
  & 95,68\%, 0,478 s/ảnh
  & Nhận dạng biển, không có nghiệp vụ bãi xe \\
Pradhan và cộng sự~\autocite{pradhan2025iotparking}
  & Raspberry Pi 4
  & 88,9\% tổng thể
  & Bãi xe IoT đầy đủ: phiên, phí, chỗ trống \\
Ammar và cộng sự~\autocite{ammar2023multistage}
  & Jetson Xavier AGX
  & 69\% (sau voting)
  & Đa khung, đa camera, không có dashboard \\
Đề tài (PC, tập vn\_plate)
  & CPU (dev)
  & 80,2\% (OCR crop), detector mAP@0,5 = 0,983
  & Đầy đủ: nhận diện, phiên, phí, dashboard; màu biển \\
Đề tài (Raspberry Pi 5)
  & Raspberry Pi 5
  & Như PC (khớp dev 50/50), 0,464 s/lượt
  & Cùng hệ thống, thêm worker tại cổng \\
\bottomrule
\end{tabular}
\end{table}

So với Tran và Bui, đề tài đọc được ký tự ``Đ'' và tích hợp toàn bộ nghiệp vụ bãi xe. So với Pradhan và cộng sự, đề tài bổ sung phân
loại kiểu dáng xe và nhận màu nền biển làm tín hiệu phân nhóm. So với Ammar và cộng sự, đề
tài nhắm phần cứng thương mại phổ thông thay vì Jetson chuyên dụng và có dashboard tra cứu
phiên đầy đủ.
